\documentclass[parskip=half, fontsize=12pt, fleqn]{scrartcl}

\usepackage[ngerman]{babel}
\usepackage{tabto} % alignment in enumerate
\usepackage[ngerman]{babel}         % Deutsches Sprachpaket
\usepackage{iftex}
\ifPDFTeX
   \usepackage[utf8]{inputenc}% Eingaben codieren
   \usepackage{fourier}
   % \usepackage[T1]{fontenc}   % Umlaute codieren, Silbentrennung
\else
   \usepackage{lmodern}
   % \usepackage[no-math]{fontspec}
   % \setmainfont{Utopia Std}
\fi

\usepackage{amsmath, amssymb}       % Mathe, \mathbb{R}
\usepackage{amsthm,amstext}         % Theoreme, \text im Mathe-Modus
\usepackage{mathtools}              % \Aboxed für Boxen in Align Umgebungen
\usepackage[arrowdel]{physics}      % Ableitungen \dv{B}{t} \pdv \dd{t}
\usepackage[left=2.5cm, right=2.5cm, top=2.8cm, bottom=2.8cm]{geometry}
\usepackage{graphicx}               % \includegraphics
\usepackage[extendedchars]{grffile} % extends file name processing of graphics
\usepackage[section]{placeins}      % \Floatbarrier
\usepackage{wrapfig}                % Bilder umfließen
\usepackage{enumerate}              % Aufzählungen
\usepackage{enumitem}               % \begin{enumerate}[label=\alph*]
\usepackage{footnote}               % Fußzeilen
\usepackage{booktabs}               % publication quality tables
\usepackage{tikz, pgfplots}         % TIKZ ist kein Zeichenprogramm
\usepackage[europeanvoltages,europeanresistors]{circuitikz}
\usepackage{bm}                     % bold symbols \bm{r}
\usepackage{dsfont}                 % identity matrix \mathds{1}
\usepackage{esint}                  % Doppelintegrale
\usepackage{mathrsfs}               % \mathscr{} statt \mathcal{}
\usepackage{placeins}               % FloatBarrier
\usepackage{subcaption}
\usepackage{multirow}
\usepackage{mdframed}               % Deckblatt
\usepackage{xcolor}                 % Deckblatt
\usepackage[headsepline]{scrlayer-scrpage}
% \usepackage{fancyhdr}               % Kopfzeile
\usepackage{aligned-overset}        % Ausrichtungen mit stackrel oder overset
\usepackage{cancel}
\usepackage{float}
\usepackage{cite}
\usepackage[version=4]{mhchem}                 % Chemistry Package
\usepackage{multicol}
\usepackage{setspace}% http://ctan.org/pkg/setspace
% \usepackage{pdfpages}             % insert whole pdf files
\usepackage[normalem]{ulem}
\usepackage{tabto}                  % alignment in enumerate
\usepackage{tcolorbox}
\tcbuselibrary{theorems}
\newtcolorbox{mymathbox}[1][]{colback=white, sharp corners, #1}

\definecolor{FSUblau}{cmyk}{1,0.7,0.1,0.5}
\definecolor{PAForange}{cmyk}{0.1,0.7,1,0}
\definecolor{Gruen}{cmyk}{1,0.1,0.7,0.5}

\usepackage[final,
    pdfauthor={Martin Beyer},
    pdffitwindow=false,     % resize document window to fit document size
    pdftoolbar=true,        % Adobe Toolbar
    bookmarks=true,         % Anzeigen der Kapitel
    bookmarksopen=true,
    bookmarksopenlevel=0,
    bookmarksnumbered=true,
    colorlinks=true,        % fuer Druckversion auf "false"
    linkcolor=FSUblau,         % Table of Contents, Footnotes
    urlcolor=FSUblau,          % fuer eingebunden URLs
    citecolor=FSUblau,         % Equations, References
    filecolor=FSUblau,
    pdfborder={0 0 0},      % keine Rahmen um Verlinkungen: {0 0 0}
    pagebackref=false
]{hyperref}

\pgfplotsset{compat=1.18}
\newcommand\mydots{\makebox[1em][c]{.\hfil.\hfil.}}
\newcommand{\minus}{\scalebox{0.75}[1.0]{$-$}}
\newcommand{\e}{\mathrm{e}}
\renewcommand{\i}{\mathrm{i}}
\newcommand{\hp}{\hphantom}

\usepackage[detect-all,
            locale=DE,
            exponent-product = \cdot,
            per-mode=fraction]{siunitx}
\usepackage[position=below,
            format=hang,
            figurename=Abb.,
            labelfont={bf},
            font=small]{caption}

\sisetup{range-phrase = {\mydots}}
% Commands
\usetikzlibrary{positioning,intersections,calc,external}
\usetikzlibrary{decorations.pathreplacing,calligraphy}
\usetikzlibrary{tikzmark}
\usepgfplotslibrary{fillbetween, groupplots}
\pgfplotsset{
tick label style={font=\small},
label style={font=\small},
legend style={font=\footnotesize},
every axis post/.style={legend cell align={left}}}
\tikzstyle{every node}=[font=\small]

\setlength{\parindent}{0px}         % keine Absätze durch Leerzeilen im Code
\numberwithin{equation}{section}

% Supress warning of incompatibility of SIunitx and physics with \qty 
\ExplSyntaxOn
\msg_redirect_name:nnn { siunitx } { physics-pkg } { none }
\ExplSyntaxOff

% Remove page number from \thispagestyle{empty}
\makeatletter\let\ps@plain\ps@fancy\makeatother

% Deckblatt
\newcommand{\HRule}{\rule{\linewidth}{0.5mm}}
\newcommand{\Deckblatt}[5][\LaTeX-Satz und Design von Martin Beyer]{
  \begin{titlepage}
    \center
    \textsc{\LARGE Friedrich-Schiller-Universität Jena\\[1ex]
    \Large Physikalisch-Astronomische-Fakultät}
    \begin{figure}[h!]
       \centering
       \includegraphics[scale=0.75]{uni-Logo_neu.pdf}
    \end{figure}\\
    \vspace{2em}
    \textsc{\Large #2}\\[0.35cm]
    \HRule \\[0.4cm]
    \begin{doublespacing}
      {\Huge\bfseries #3}\\[0.15cm]
      \end{doublespacing}
    \HRule \\[0.5cm]
    \textsc{\Large #4}\\[0.5cm]
    {\large Version: \today}\\
    \vfill
    \begin{mdframed}[backgroundcolor=gray!20]
      \begin{center}
        #1
      \end{center}
    \end{mdframed}
  \end{titlepage}

  % Set the default page style to scrheadings
  \pagestyle{scrheadings}
  \automark{subsection}
  \renewcommand{\headfont}{\upshape}

  % Configure headers and footers
  \clearpairofpagestyles
  \ihead{\rightmark} % inner header
  \ohead{\textbf{#5}} % outer header
  \ofoot*{\bfseries\thepage} % outer footer
  % \pagestyle{fancy}
  % \fancyhead[R]{\textbf{#5}}
  % \fancyfoot[R]{\bfseries\thepage}
  % \fancyfoot[C]{}
  % \fancyhead[L]{\rightmark}

  % \fancypagestyle{plain}{
  %   \fancyfoot[R]{\bfseries\thepage}
  %   \fancyfoot[C]{}
  %   \fancyhead[R]{}
  %   \fancyhead[L]{}
  %   \renewcommand{\headrulewidth}{0pt}
}
% \renewcommand{\sectionmark}[1]{\markright{#1}}
% \renewcommand{\headrulewidth}{0.5pt}
% \renewcommand{\footrulewidth}{0pt}




\makeatletter
\renewcommand\tagform@[1]{}
\makeatother

\setlength{\mathindent}{0cm}

\newgeometry{left=2cm, right=2cm, top=2cm, bottom=2.5cm}
\setenumerate{labelindent=1em,labelsep=0.5cm,leftmargin=*}
\usetikzlibrary{tikzmark}
\everymath{\displaystyle}
\newcommand{\Semester}{2023/24}
\newcommand{\hp}{\hphantom}

\newcommand{\Titelbanner}[2]{
    \begin{center}
        \textsc{
        \LARGE Auffrischungskurs Mathematik\\[0.5cm]
        \large  -- Ausführliche Lösungen für Tutoren --\\[0.5cm]}
        \footnotesize WS \Semester
    \end{center}
    \paragraph{Thema #1:} \hspace{0.2cm}
    \begin{minipage}[t]{0.8\linewidth}
        #2
    \end{minipage}
    \vspace{0.7cm}
}

%%%% 
% \includeonly{Lösungen_Lang/Thema_05} 
%%%%

\begin{document}
\graphicspath{{./Bilder/}{../}}
% \Deckblatt{Wintersemester \Semester}
%           {Mathematik - Ein Vorkurs für Studienanfänger}
%           {Martin Beyer}
%           {Mathematik}

\pagestyle{plain}
\renewcommand{\headrulewidth}{0pt}
\renewcommand{\footrulewidth}{0pt}
\pagenumbering{gobble}
% \tableofcontents
% \newpage
  

\Titelbanner{1}{Grundrechenarten\\
                Brüche\\
                Potenzen\\
                Wurzeln}

\begin{minipage}[t]{0.6\linewidth}
    \textbf{Aufgabe 1: } \emph{Bruchrechnung}
    \begin{enumerate}[label=(\alph*)]
        \item $b-a$
        \item $\dfrac{a}{b}$
        \item $\dfrac{x}{y}-1$
        \item $\dfrac{n^2}{n-1}$
        \item $\dfrac{x+y}{x-y}$
        \item $1$
        \item $\dfrac{a^2}{\sqrt{2}}$
    \end{enumerate}
\end{minipage}
\hfill
\begin{minipage}[t]{0.37\linewidth}
    \textbf{Aufgabe 2: } \emph{Potenzgesetze}
    \begin{enumerate}[label=(\alph*)]
        \item $\left( \dfrac{a+b}{x-y} \right) ^n$
        \item $a^xb^yc^z$
        \item $(a+b)a^2b^{n-2}$
        \item $(a-1)a^{n-2}$
        \item $\dfrac{x^5}{a^{14}b^{19}y^{12}}$
    \end{enumerate}
\end{minipage}\\[2cm]
%
\begin{minipage}[t]{0.6\linewidth}
\textbf{Aufgabe 3: } \emph{Umformungen mit Wurzelausdrücken}
    \begin{enumerate}[label=(\alph*)]
        \item $\dfrac{1+ab}{1-ab}$
        \item $\dfrac{1}{m}$
        \item $a\left(\sqrt[4]{ab}+\sqrt{b}\right)$
    \end{enumerate}
\end{minipage}
\hfill
%
\begin{minipage}[t]{0.4\linewidth}
\textbf{Aufgabe 4: } \emph{Algebraische Umformungen}
    \begin{enumerate}[label=(\alph*)]
        \item $x=\dfrac{1}{(a-b)(m+n)}$
        \item $x=\dfrac{a-b}{a+b}$
        \item $x=\dfrac{3}{4}$
        \item $x=(a+b)^2$
        \item $x=\left|y^2-1\right|$
    \end{enumerate}
\end{minipage}
\Titelbanner{2}{Lineare Gleichungssysteme}

\paragraph{Tafelbeispiel:} $\displaystyle \frac{10(x+y) + 3}{x-2y+4} = 1, \quad \frac{36x-3y}{7(x-y)+3} = 3$
\begin{enumerate}
    \item Sortieren nach $x$ und $y$ 
    \begin{align*}
        9x +12 y &= 1 \quad (1) \\
        15x+18y &= 9 \quad (2)
    \end{align*}
    \item Prüfe auf lineare Unabhängigkeit, d.\,h. Berechne $a_1 b_2 - a_2 b_1$ (Determinante, vgl. Vorlesung): 
    \begin{align*}       
        9\cdot 18 - 15\cdot 12 = 162-180 \neq 0 \quad \Rightarrow \qq{genau ein Lösungspaar ($x,y$)}
    \end{align*}      
    \item Lösen, bspw. durch geschickte Linearkombination (Elimination von $y$)
    \begin{align*}
        2\cdot (2) - 3\cdot (1): \quad \hp{1}3x &= \hp{-}15 \quad \Rightarrow x=5 \\
        \qq{in} (1): \quad 12y &= -44 \quad \Rightarrow y = -\frac{11}{3} \\
        \Rightarrow \qq{Lösungsmenge} \mathbb{L} &= \uuline{\Bigg\{\qty(5;-\frac{11}{3})\Bigg\}}
    \end{align*}
\end{enumerate}

\paragraph{Aufgabe 1: } \emph{Zwei lineare Gleichungen mit zwei Unbekannten} \hfill Ziel (a) bis (d)\\[0.2cm]
\begin{enumerate}[label=(\alph*)]
    \item $~$\\[-1.45 cm]\begin{align*}
        &33x+12y=25 \quad (1) \\ 
        &11x-\hp{1}3y=\hp{2}6 \quad (2) \\
        &(1) - 3\cdot (2): \quad 21y = 7 \quad\Rightarrow \quad y = \frac{1}{3}, \qq{in (2):} x = \frac{7}{11} \qquad \Rightarrow \qquad \uuline{\mathbb{L} = \Bigg\{\qty(\frac{7}{11};\frac{1}{3})\Bigg\}}
    \end{align*}
    \item $~$\\[-1.45cm] 
    \begin{align*}
        &\begin{rcases}
            \frac{2x + 3y}{3x-y} = \frac{17}{9} \\
            \dfrac{3x+4y}{6x-1}=2
        \end{rcases} \quad \Longrightarrow \quad 
        \begin{array}{l}
            3x-4y = 0 \quad (1) \\
            9x -4y = 2 \quad (2)
        \end{array} \\
        &\hp{3\cdot}\;\,(1) -(2): \quad 6x = 2 \quad\Rightarrow\quad x = \frac{1}{3}\\
        &3\cdot(1) -(2): \quad 8y = 2 \quad\Rightarrow\quad y = \frac{1}{4} \qquad \Rightarrow \qquad \uuline{\mathbb{L} = \Bigg\{\qty(\frac{1}{3};\frac{1}{4})\Bigg\}}
    \end{align*}
    \item $~$\\[-1.45cm] 
    \begin{align*}
        &\begin{rcases}
            \frac{x+2}{y+3}=\frac{1}{3} \\
            \frac{y+3}{2y-5x}=\frac{3}{5}
        \end{rcases} \quad \Longrightarrow \quad 
        \begin{array}{l}
            \hp{1}3x-y = \;\;\;-3 \quad (1) \\
            15x -y = -15 \quad (2)
        \end{array} \\
        &(1) - (2): \quad -12x = 12 \quad \Rightarrow \quad  x = -1\\
        &\quad\,\text{in } (1): \quad \hp{-12x = 12\quad \Rightarrow \quad} \,y = 0 \qquad \Longrightarrow \qquad \uuline{\mathbb{L} = \big\{\qty(-1;0)\big\}}
    \end{align*}
    \item $~$\\[-1.45cm] 
    \begin{align*}
        &ax+by=2a \quad (1)\\
        &\frac{x}{b}-\frac{y}{a}=\frac{2}{a} \quad (2)\\
        &(1) - ab \cdot (2): \;\;\,\cancel{2}by = \cancel{2}(a-b) \quad \Rightarrow \quad y = \frac{a}{b}-1 \\
        &\hp{(1)-}\text{in } (1): \quad ax + \cancel{a} -b = \cancel{2}a \quad \Rightarrow \quad x=\frac{b}{a}+ 1 \qquad \Longrightarrow \qquad \uuline{\mathbb{L} = \Bigg\{\qty(\frac{b}{a}+1;\frac{a}{b}-1)\Bigg\}}
    \end{align*}
    \item $~$\\[-1.45cm] 
    \begin{alignat*}{2}
        &x+14y=\frac{1}{\sqrt{2}}-7\sqrt{2} \;\;\;\quad (1)\\
        &3\sqrt{2}x-\frac{y}{\sqrt{3}}=3+\frac{1}{\sqrt{6}} \quad (2)\\
        &(1) - \frac{1}{3\sqrt{2}}\cdot (2): \quad\cancel{\qty(14+\frac{1}{3\sqrt{6}})}y &&= \bcancel{\frac{1}{\sqrt{2}}} - 7\sqrt{2} - \bcancel{\frac{1}{\sqrt{2}}} - \frac{1}{3\sqrt{12}} \\
        & &&= -\frac{1}{\sqrt{2}} \cancel{\qty(14+\frac{1}{3\sqrt{6}})} \quad \Rightarrow \quad y = -\frac{1}{\sqrt{2}} \\
        & \text{in } (1): \quad x-\cancel{7\sqrt{2}} = \frac{1}{\sqrt{2}} - \cancel{7\sqrt{2}} \quad &&\Rightarrow \quad x = \frac{1}{\sqrt{2}} \qquad \Longrightarrow \qquad \uuline{\mathbb{L} = \Bigg\{\qty(\frac{1}{\sqrt{2}};-\frac{1}{\sqrt{2}})\Bigg\}}
    \end{alignat*}
    \item $~$\\[-1.45cm] 
    \begin{alignat*}{3}
        &\frac{x}{a+b}+\frac{y}{a-b}=a+b \quad (1)\\
        &\frac{x}{a}-\frac{y}{b}=2b \hspace{1.68cm}\quad (2) \\
        &(a+b) \cdot (1) - a\cdot(2): \quad \qty(\frac{a+b}{a-b}+\frac{a}{b}) y &&= a^2 +b^2 \\
        && \hspace{-3.5cm}\bcancel{\qty(\cancel{ab}+b^2 + a^2 - \cancel{ab})}y &= b(a-b)(a^2 +b^2) \\
        && \hspace{-1cm}\Rightarrow y &= b(a-b) \\
        &\text{in }(1): \quad \frac{x}{a+b} +\cancel{b} = a+ \cancel{b} & \hspace{-1cm}\Rightarrow x &= a(a+b) \quad \Longrightarrow \quad \uuline{\mathbb{L} = \big\{\qty(a(a+b);b(a-b))\big\}}
    \end{alignat*}
    \item $~$\\[-1.45cm] 
    \begin{alignat*}{3}
        &39x-38y=1 \quad (1) \\
        &91x-57y=4 \quad (2) \\
        &2\cdot(2)-3\cdot(1): \quad (182-117)x = 5 \quad &&\Rightarrow \quad x = \frac{1}{13} \\
        &\text{in } (1): \quad 38y = 2 &&  \Rightarrow \quad y = \frac{1}{19} \qquad \Longrightarrow \qquad \uuline{\mathbb{L} = \Bigg\{\qty(\frac{1}{13};\frac{1}{19})\Bigg\}}
    \end{alignat*}
\end{enumerate}

\newpage
\paragraph{Aufgabe 2: } \emph{Drei Gleichungen mit drei Unbekannten} \hfill Ziel: (a) und (b)\\[0.2cm]
%
\emph{Tafelbeispiel: }
\begin{alignat*}{2}
    x+ \hp{2}y + \hp{4}z &= 9 &&\quad (1) \\
    x+2y+4z&=15 &&\quad (2) \\
    x+3x+9z&=23 &&\quad (3)
\end{alignat*}
Vorgehen: Wir wollen eine Linearkombination der drei Gleichungen finden, sodass (bspw.) $x$ und $y$ verschwinden; d.\,h. in $a\cdot(1)+b\cdot(2)+c\cdot(3)$ soll $a+b+c =0$ sein und $a+2b+3c=0$. Wir haben offenbar die Freiheit, $a=1$ zu wählen, also $b=-2$ und $c=1$.
\begin{align*}
    &(1)-2\cdot(2)+(3): \quad 2z=2 \quad \Rightarrow \quad z=1 \\
    &\text{in }(1): \quad x+\hp{2}y = 8\;\;\; \quad(4)\\
    &\text{in }(2): \quad x+2y = 11 \quad (5)
\end{align*}
Die Gleichungen (4),(5) können nun mit den früheren Methoden gelöst werden. \\
Optionale Kreativlösung: Angenommen $x,y \in \mathbb{N}$. Nach (5) muss $x$ ungerade sein, damit nach (4) auch $y$. Die einzigen positiven Zerlegungen der Zahl 8 in zwei ungerade Zahlen sind (1,7) und (3,5). Nur letztere löst (4) und (5). Wir erhalten $\uuline{\mathbb{L} = \{(5;3;1)\}}$.

\emph{Lösung:}

\begin{enumerate}[label=(\alph*)]
    \item $~$\\[-1.45 cm]
    \begin{align*}
        &\hp{3}x-\hp{5}y+5z=5 \quad (1)\\
        &3x+7y-5z=5\quad (2)\\
        &\hp{3}x+\hp{5}y-\hp{5}z=1 \quad (3) \\
        &2\cdot(1) + (2) - 5\cdot(3): \quad 10z=10 \quad \Rightarrow \quad z=1 \\
        &\begin{rcases}
        \text{in }(1): \quad x-y = 0\\
        \text{in }(3): \quad x+y = 2
        \end{rcases} \quad x=y=1 \qquad \Longrightarrow \qquad \uuline{\mathbb{L} = \big\{\qty(1;1;1)\big\}}
    \end{align*}
    \item $~$\\[-1.45cm] 
    \begin{align*}
        &3x-4y+3z=\hp{-}4\quad (1)\\
        &\minus x+\hp{4}y-\hp{5}z=-2\quad (2)\\
        &7x+4y-5z=\hp{-}0 \quad (3) \\
        &(1) + 8\cdot (2) -(3): \quad -12x = -12 \quad \Rightarrow \quad x=1 \\
        &\begin{rcases}
            \text{in }(1): \quad -4y+3z = \hp{-}1\\
            \text{in }(2): \quad \hp{-4}y-z = -1
            \end{rcases} \quad y=2, z=3 \qquad \Longrightarrow \qquad \uuline{\mathbb{L} = \big\{\qty(1;2;3)\big\}}
    \end{align*}
    \item $~$\\[-1.45cm] 
    \begin{align*}
        &x+y=b+a\quad (1)\\
        &x+z=a+c\quad (2) \qquad \text{Lösung durch Draufschauen.}\\
        &y+z=c+b\quad (3)\\
        &(1)+(2)-(3): \quad 2x=2a \\
        &(1)-(2)+(3): \quad 2y=2b \\
        &(2)+(3)-(1): \quad 2z=2c \qquad \Longrightarrow \qquad \uuline{\mathbb{L} = \big\{\qty(a;b;c)\big\}}
    \end{align*}
    \newpage
    \item $~$\\[-1.45cm] 
    \begin{align*}
        &\hp{\minus}6x-4y+\;\;\;\;8z=0\quad (1) \\
        &\minus2x+\hp{4}y-\,\hp{11}z=0\quad (2) \qquad \text{Da $(1)-3\cdot(2)=(3)$, ist das System unterbestimmt.}\\
        &12x-7y+11z=0\quad(3)\\
        &(1)+3\cdot(2): \quad -\hp{2}y+5z = 0 \quad \Rightarrow y=5z \\
        &(1)+4\cdot(2): \quad -2x+4z=0 \quad \Rightarrow x=2z \qquad \Longrightarrow \qquad \uuline{\mathbb{L} = \big\{\qty(2z;5z;z)\mid z\in\mathbb{R}\big\}
        }
    \end{align*}
\end{enumerate}

% 
\textbf{Aufgabe 3: } \emph{Parametrisierung von Lösungsmengen}\\[0.2cm]
Geben Sie die Lösungsmenge der Gleichung $13x-7y=1$ an für
\begin{align*}
&\text{(a) } \hspace{0.2cm} x,y\in\mathbb{R}\,;\\[0.2cm]
&\text{(b) } \hspace{0.2cm} x,y\in\mathbb{N}\,.
\end{align*}
(a) Für $x,y \in \mathbb{R}$ ist nichts zu beachten. Wir Lösen die Gleichung nach $x$ auf und erhalten 
\begin{align*}
    x = \frac{1+7y}{13} \qquad \Longrightarrow \qquad \uuline{\mathbb{L} = \Bigg\{\qty(\frac{1+7\lambda}{13};\lambda)\mid \lambda \in\mathbb{R}\Bigg\}}.
\end{align*}
(b) Für $x,y \in \mathbb{N}$ muss $1+7n, n\in\mathbb{N}$, ein Vielfaches von 13 sein. Die kleinste natürliche Zahl, für die das gilt, ist $n=11$ (ausprobieren), d.\,h. alle übrigen Teiler liegen in 13er-Schritten darüber.
\begin{table}[htp]
    \centering
    \begin{tabular}{l r r r r r r r r r r r r}
        \toprule 
        $n$ & 1 & 2 & 3 & 4 & 5 & 6 & 7 & 8 & 9 & 10 & 11\\
        \midrule
        $7n+1$ & 8 & 15 & 22 & 29 & 36 & 43 & 50 & 57 & 64 & 71 & 78 \\
        13-Reihe & 13 & 26 & 39 & 52 & 65 & 78 
    \end{tabular}
\end{table}

Setzen wir also $\lambda = 11+13n$ in $\mathbb{L}$ von oben ein, ergibt sich $\uuline{\mathbb{L} = \Bigg\{\qty(6+7n;11+13n)\mid n \in\mathbb{N}_0\Bigg\}}$.

\paragraph{Aufgabe 4: } \emph{Gleichungssysteme} \hfill Ziel (a) und (b)\\[0.2cm]

\emph{Tafelbeispiel:} In dieser Aufgabe werden nicht linear aussehende Gleichungen auf lineare Gleichungssysteme zurückgeführt. 
\begin{alignat*}{2}
    (1-x)^2 &= 5y^2 - 4(x-2)^2 \quad &&(1) \\
    2(x^2-y^2) &= 6x-5 &&(2).
\end{alignat*}
Sortieren wir das Gleichungssystem um, ergibt sich 
\begin{alignat*}{2}
    &5(x^2-y^2)-18x + 17 = 0 \quad (3)\\
    &2(x^2-y^2)-\hp{1}6x + \hp{1}5 = 0 \quad (4)\\
    &(3)-\frac{5}{2}\cdot(2): \quad -3x + \frac{9}{2} = 0 \quad\Rightarrow \quad &&x=\frac{3}{2} \\
    &\text{in }(2): \quad y^2 =x^2 -3x + \frac{5}{2} = \frac{1}{4} \quad &&y= \pm \frac{1}{2} \qquad \Longrightarrow \qquad \uuline{\mathbb{L} = \Bigg\{\qty(\frac{1}{2};\pm \frac{1}{2})\Bigg\}}. 
\end{alignat*}

\emph{Lösung:}

\begin{enumerate}[label=(\alph*)]
    \item $~$\\[-1.45 cm]
    \begin{align*}
        &x^2+y^2=2(xy+2) \quad(1) \quad \Rightarrow \quad (x-y)^2 = 4 \Rightarrow x-y = \pm 2 \quad (1a)\\
        &x\hp{^2}+y\hp{^2} =6 \hp{(xy+2)} \quad(2)\\
        &(1a)+(2): \quad 2x = \begin{cases}
            8, & \text{oberes VZ} \\
            4, & \text{unteres VZ}
        \end{cases}\\
        &\text{in }(2): \qquad\quad\; y = \begin{cases}
            2, & \text{oberes VZ} \\
            4, & \text{unteres VZ}
        \end{cases} \qquad \Longrightarrow \qquad \uuline{\mathbb{L} = \Bigg\{\qty(4;2),(2;4)\Bigg\}}
    \end{align*}
    \item $~$\\[-1.45cm]
    \begin{align*}
        &\frac{12}{\sqrt{x-1}}+\frac{5}{\sqrt{y+\frac{1}{4}}}=5\quad (1)\\ 
        &\frac{8}{\sqrt{x-1}}+\frac{10}{\sqrt{y+\frac{1}{4}}}=6\quad(2)\\
        &2\cdot(1)-(2):\hp{3\cdot}\; \quad 16 = 4\sqrt{x-1}\;\;\; \quad \Rightarrow \quad x=17 \\
        &2\cdot(1)-3\cdot(2): \quad 20 = 8\sqrt{y+\frac{1}{4}} \quad \Rightarrow \quad y=6 \qquad \Longrightarrow \qquad \uuline{\mathbb{L} = \Bigg\{\qty(17;6)\Bigg\}}
    \end{align*} 
    Optionaler Hinweis: nicht immer muss alles mühsam quadriert und aufgelöst werden, sondern man kann auch versuchen draufzuschauen. 
    \begin{itemize}
        \item Was ergibt 8 multipliziert mit 20? Antwort: $\textstyle\frac{10}{4}=\frac{5}{2}$
        \item Die Wurzel woraus ist $\textstyle\frac{5}{2}$? Antwort: $\textstyle\frac{25}{4}$. 
        \item Damit ist $y = \frac{25}{4}-\frac{1}{4}= \frac{24}{4} = 6$.
    \end{itemize}
    \vspace{3mm}
    \item $~$\\[-1.45cm]
    \begin{align*}
        &\frac{x^2-y^2}{2x+3}+y^2=(x+y)x-xy \quad \Rightarrow \quad \frac{x^2 -y^2}{2x+3} = x^2 - y^2\quad(1) \\
        &y-2x=3 \quad (2)
    \end{align*}
    \begin{align*}
        \text{Fall } x^2 \neq y^2 \quad
        \begin{rcases}
            (1): 1=2x+3 \quad \Rightarrow x=-1 \\
            \text{in }(2): \quad y=1
        \end{rcases} \text{ Widerspruch}
    \end{align*}
    \begin{align*}
        \text{Fall } x = \pm y \quad (2): \quad y=
        \begin{cases}
            -3, &\text{oberes VZ}\\
            1, &\text{unteres VZ}
        \end{cases} \qquad \Longrightarrow \qquad \uuline{\mathbb{L} = \Bigg\{\qty(-3;-3),(-1;1)\Bigg\}}
    \end{align*}
\end{enumerate}

\newpage
\paragraph{Aufgabe 5: } \emph{Ungleichungssysteme}
\begin{enumerate}[label=(\alph*), labelindent=1em,labelsep=0.5cm]
\item Es ist die Lösungsmenge des folgenden Systems an Ungleichungen zu skizzieren. Welche der Ungleichungen können weggelassen werden, ohne dass sich die Lösungsmenge ändert?

\begin{minipage}{0.2\textwidth}
    \begin{alignat*}{3}
        2x-3y&\geq -6\,,\\
        x-2y&< 11\,,\\
        x&>-y-1\,,\\
        x&<5\,,\\
        x&\geq 0\,,\\
        y&\geq 0
    \end{alignat*}
\end{minipage}
\begin{minipage}{0.8\textwidth}
    \centering
        \begin{tikzpicture}
            \begin{axis}[disabledatascaling, axis lines=middle, xlabel={x}, ylabel={y}, height=7.8cm, domain=-2:10, xmin=-3,xmax=11, ymin=-8, ymax=9, axis equal, legend pos= outer north east]
                \addplot[no marks, black, thick, ]{2/3*x+2};
                \addplot[no marks, FSUblau, thick]{0.5*x-11/2};
                \addplot[no marks, Gruen, thick]{-x-1};
                \draw[thick, PAForange, thick] (5,-8)--(5,8)node[right]{$x<5$};
                \fill[FSUblau, fill opacity =0.3] (0,0)--(5,0) --(5,16/3) -- (0,2) -- cycle;
                \legend{$\textstyle y = \frac{2}{3}x+2$, $\textstyle y = \frac{1}{2}x-\frac{11}{2}$, $\textstyle y = -x-1$};
            \end{axis}
        \end{tikzpicture}
\end{minipage}
Offenbar können die zweite und dritte Gleichung weggelassen werden, ohne dass sich etwas am schraffierten Gebiet ändert. 

\item \textbf{*} Welches Gebiet im ersten Oktanden ($x\ge 0$, $y\ge 0$, $z\ge 0$) wird durch die folgenden Ungleichungen definiert?
\begin{align*}
    x+y&\ge z\,,\\
    x+z&\ge y\,,\\
    y+z&\ge x
\end{align*}

\begin{minipage}{0.4\textwidth}
    Die Ungleichungen beschreiben eine dreiseitige (unendlich große) Pyramide, deren Spitze auf dem Koordinatenursprung liegt und deren Kanten jeweils die $x-y-$, $x-z-$ und $x-z$-Ebene halbieren.
\end{minipage}
\begin{minipage}{0.6\textwidth}
    \centering
        \begin{tikzpicture}
            \draw[thick, -{latex}](0,0) -- (-30:2)node[right]{$y$};
            \draw[thick, -{latex}](0,0) -- (90:2)node[above]{$z$};
            \draw[thick, -{latex}](0,0) -- (210:2)node[left]{$x$};
            \foreach \x in {0.25,0.5,...,2.5}{
                \draw[thick] (-90:\x) -- (30:\x) -- (150:\x) -- cycle;
                \fill[fill opacity=0.01] (0,0) -- (-90:2.5) -- (150:2.5) -- cycle;
                \fill[fill opacity=0.03] (0,0) -- (-90:2.5) -- (30:2.5) -- cycle;
            }
        \end{tikzpicture}
\end{minipage}

\end{enumerate} 
\Titelbanner{3}{Quadratische Gleichungen und Gleichungssysteme}

\paragraph{Aufgabe 1: } \emph{Quadratische Gleichungen} \hfill Ziel: (a) und (b)\\[0.2cm]
Lösen Sie die folgenden Gleichungen jeweils für $x$ durch quadratische Ergänzung und kontrollieren Sie das Ergebnis mit der $pq$-Formel.
\begin{enumerate}[label=(\alph*)]
    \item $x^2-10x+9=0$
    \item $x^2+x-12=0$
    \item $x^2-\sqrt{8}x+1=0$
\end{enumerate}

\paragraph{Aufgabe 2: } \emph{Wurzeln quadratischer Gleichungen}
\begin{enumerate}[label=(\alph*)]
    \item Stellen Sie $\textstyle\frac{a}{b}-\frac{b}{a}$ als Produkt zweier Faktoren dar, deren Summe gleich $\textstyle\frac{a}{b}+\frac{b}{a}$ ist.
    \item Bestimmen Sie in der Gleichung $5x^2-kx+1=0$ den Koeffizienten $k$ so, dass die Differenz der Wurzeln 1 ergibt.
    \item Wählen Sie die Koeffizienten der quadratischen Gleichung $x^2+px+q=0$ so, dass die Wurzeln der Gleichung gleich $p$ und $q$ sind.
    \item Gegeben ist die quadratische Gleichung $ax^2+bx+c=0$. Gesucht ist diejenige neue quadratische Gleichung, deren Wurzeln gleich
    \begin{itemize}[labelindent=1em,labelsep=0.5cm]
        \item dem Doppelten der Wurzeln der gegebenen Gleichung,
        \item den reziproken Werten der Wurzeln der gegebenen Gleichung
    \end{itemize}
    sind.
\end{enumerate}

\paragraph{Aufgabe 3: } \emph{Gleichungssysteme}\\[0.2cm]
Lösen Sie die folgenden Gleichungssysteme jeweils für $x$ und $y$.\\[0.2cm]
\begin{minipage}[t]{0.25\linewidth}
(a)\vspace{-2.6em}
\begin{align*}
x+y^2&=7\,,\\ xy^2&=12
\end{align*}
\end{minipage}\hspace{0.05\linewidth}
\begin{minipage}[t]{0.3\linewidth}
(b)\vspace{-2.6em}
\begin{align*}
x+xy+y&=11\,,\\ x^2y+xy^2&=30
\end{align*}
\end{minipage}\hspace{0.05\linewidth}
\begin{minipage}[t]{0.3\linewidth}
    (c)\vspace{-2.8em}
\begin{align*}
x^2+y^2&=\frac{5}{2}xy\,,\\ x-y&=\frac{1}{4}xy
\end{align*}
\end{minipage}\\
\emph{Hinweis:} Substitution $z_1=xy$, $z_2=x+y$ in Aufgabe (b).

\paragraph{Aufgabe 4: } \emph{Wurzelgleichungen} \hfill Ziel: (a) bis (c)\\[0.2cm]
Lösen Sie die folgenden Gleichungen jeweils für $x$.
\begin{enumerate}[label=(\alph*)]
    \item $\sqrt{3x+1}-\sqrt{x-1}=2$
    \item $\sqrt{x+a}=a-\sqrt{x}$
    \item $\sqrt{a^2-x}+\sqrt{b^2-x}=a+b$
    \item $\sqrt{x+1+\sqrt{3x+4}}=3$
    \item $\sqrt{2x-1}+\sqrt{x-\num{1,5}}=\dfrac{6}{\sqrt{2x-1}}$
\end{enumerate}
%
\paragraph{Aufgabe 5: } \emph{Nullstellensuche} \hfill Ziel: (a) bis (c)\\[0.2cm]
Die folgenden Terme sind als Produkte von Linearfaktoren darzustellen.
\begin{enumerate}[label=(\alph*)]
    \item $x^2+2x-15$
    \item $4x^2+8x-5$
    \item $ax^3+bx^2+adx^2+bdx$
    \item $(a-x)^2+(x-b)^2-a^2-b^2$
    \item $a\sqrt{8}x^2-2kax-3ak+ax\sqrt{18}$
\end{enumerate}
\Titelbanner{4}{Umgang mit Polynomen höheren Grades\\
                Das Summenzeichen}

\paragraph{Aufgabe 1: } \emph{Nullstellensuche} \hfill Ziel: (a) bis (c)\\[-4mm]
\begin{enumerate}[label=(\alph*), labelindent=1em,labelsep=0.5cm]
    \item \emph{Stellen Sie das Polynom dritten Grades auf, das Wurzeln $a$, $b$ und $c$ hat.}
    \begin{align}
        f_3(x) = (x-a)(x-b)(x-c) = \uuline{x^3-(a+b+c) x^2 + (ab+ac+bc) x - abc}
    \end{align}
    \emph{Bemerkung}: Hier kann man versuchen, sich klarzumachen, wie der verallgemeinerte Satz von Vieta aussieht.
    \item \emph{Zerlegen Sie das Polynom $f_4(x)=x^3+2x^4+4x^2+2+x$ in Faktoren. Welche Aussage können Sie über dessen Nullstellen treffen?}
    \begin{align}
        f_4(x) = x^3 + \underbrace{2x^4 +4x^2 +2}_{\mathclap{2x^4+2x^2+2x^2 +2 = 2x^2(x^2+1) + 2(x^2+1)}} +x = (x^2+1)(2x^2+2) + \underbrace{x^3+x}_{x(x^2+1)} = \uuline{(x^2+1)(2x^2+x+2)}
    \end{align}
    Da weder $x^2 +1 = 0$ noch $\textstyle x^2 +\frac{1}{x}x+1=0$ eine (reelle) Lösung hat, besitzt das Polynom keine (reellen) Nullstellen.
    \item \emph{Bestimmen Sie alle Nullstellen des Polynoms $f_5(x)=x^5-3x^3+2x$.}
    Offensichtlicherweise gilt $x_1 = 0$ und wir können das Polynom mit Division durch $x$ vereinfachen. Mit der Substitution $z \equiv x^2$ erhalten wir 
    \begin{align}
        0 = z^2 -3z +2 \quad \Rightarrow \quad z_{1/2} = \frac{3}{2}\pm \sqrt{\frac{9}{4}-2} = \begin{cases}
            2 \\ 1
        \end{cases} \quad \Rightarrow \quad \uuline{x_{2/3} = \pm \sqrt{2}}, \quad \uuline{x_{4/5} = \pm 1}.
    \end{align}
    \item \emph{Bestimmen Sie die kleinste positive Nullstelle des Polynoms $\textstyle f_4(x)=1-\frac{x^2}{2}+\frac{x^4}{24}$. Setzen Sie näherungsweise $\textstyle \sqrt{3}\approx 3-\frac{\pi^2}{8}$.}
    Wir führen wieder die Substitution $z \equiv x^2$ aus und erhalten 
    \begin{align}
        0 = z^2 - 12 z + 24 \quad \Rightarrow \quad z_{1/2}= 6\pm2\sqrt{3} \qq{bzw.} x_{1/2/3/4} = \pm \sqrt{6\pm 2\sqrt{3}}.
    \end{align}
    Wir suchen nun das kleinste positive $x$, also 1.VZ: ``$+$'' und 2.VZ: ``$-$'' 
    \begin{align}
            x_0 = \sqrt{6-2\sqrt{3}} \overset{\sqrt{3}\approx 3-\frac{\pi^2}{8}}{\approx} \sqrt{6-\qty(3-\frac{\pi^2}{8})} = \uuline{\frac{\pi}{2}}.
    \end{align}
    \emph{Bemerkung}: Die Funktion $f_4(x)$ ist das dritte Taylorpolynom der Kosinusfunktion. Es gilt $\cos(x) = f_4(x) +\order{x^6}$.
\end{enumerate}

\paragraph{Aufgabe 2: } \emph{Polynomdivision}\\[0.2cm]
Berechnen Sie die folgenden Ausdrücke. Zusatz: Für welche Werte von $n$ bleibt die Polynomdivision in (d) ohne Rest?
\begin{enumerate}[label=(\alph*), labelindent=1em,labelsep=0.5cm]
    \item$~$\\[-1.3cm]
    \begin{align}
        \begin{array}{r@{} r@{} r@{} r@{} r}
            (21a^3 &{}+34a^2b\hp{)}&{}&{}+25b^3) &\;:(7a+5b) = \uuline{3a^2-7ab+5b^2} \\
          -(21a^3 &{}+15a^2b) \\ 
          \cmidrule{1-3}
                & -49a^2 b\hp{)} &{}&{}+25b^3\hp{)}\\
                &-(-49a^2 b\hp{)}&{}-35ab^2&{}\hp{+25b^3}) \\
          \cmidrule{2-4}
                & &{}35ab^2&{}+25b^3\hp{)}\\
                & &-(35ab^2&{}+25b^3)\\
          \cmidrule{3-4} 
                & & & 0
        \end{array}
    \end{align}
    \item$~$\\[-1.3cm]
    \begin{align}
        \begin{array}{r@{} r@{} r@{} r@{} r}
            (9x^3 &{}&{}-7xy^2\hp{)}&{}+2y^3) &\;:(3x-2y) = \uuline{3x^2+2xy-y^2} \\
          -(9x^3 &{}-6x^2y&{}) \\ 
          \cmidrule{1-3}
                & 6x^2y &{} -7xy^2 &{}+2y^3\hp{)}\\
                &-(6x^2y&{} -4xy^2&{}\hp{+2y^3}) \\
          \cmidrule{2-4}
                & &{}-3xy^2&{}+2y^3\hp{)}\\
                & -&(-3xy^2&{}+2y^3)\\
          \cmidrule{3-4} 
                & & & 0
        \end{array}
    \end{align}
    \item$~$\\[-1.3cm]
    \begin{align}
        \begin{array}{r@{} r@{} r@{} r@{} r@{} r}
            (25x^4 &{}&{}-a^2x^2\hp{)}&{}&{}+25a^4) &\;:(5x^2 
            + 7ax +5a^2) = \uuline{5x^2 -7ax +5a^2} \\
          -(-25x^4 &{}+35ax^3&{}+25a^2x^2) \\ 
          \cmidrule{1-3}
                & -35ax^3 &{} -24a^2x^2 &{}&{}+25a^4\hp{)}\\
                &-(-35ax^3&{} -49a^2x^2&{}-35a^3 x&{}\hp{+25a^4}) \\
          \cmidrule{2-5}
                & &{}25a^2 x^2&{}+35a^3x&{}+{25a^4}\hp{)}\\
                & &-(25a^2 x^2&{}+35a^3x&{}+{25a^4})\\
          \cmidrule{3-5} 
                & & & 0
        \end{array}
    \end{align}
    \item$~$\\[-1.3cm]
    \begin{align}
        \begin{array}{r@{} r@{} r@{} r}
            (x^2 &{} +2x \hp{)}&{}-15) &\;:(x+n) = \uuline{x+2-n + \frac{n(n-2)-15}{x+n}} \\[-3mm]
          -(x^2 &{}+nx) \\ 
          \cmidrule{1-2}
                & (2-n)x \hp{)}&{} -15\hp{)}\\
                &-((2-n)x\hp{)}&{} +n(2-n)) \\
          \cmidrule{2-3}
                & -15\hp{)}&{}+n(n-2) & \Rightarrow 0 \overset{!}{=} n^2 - 2n -15 \quad \Rightarrow \quad n_{1/2} = 1 \pm \sqrt{16}
        \end{array}
    \end{align}
    Wir erhalten damit die beiden Werte $n_{1/2}$ die den Wurzeln des quadratischen Gleichungssystems entsprechen 
    \begin{align}
        x^2 + 2x -15 = (x+5)(x-3).
    \end{align}
\end{enumerate}

\newpage
\paragraph{Aufgabe 3: } \emph{Kubische Gleichungen} \hfill Ziel: (a)\\[-4mm]
\begin{enumerate}[label=(\alph*)]
    \item \emph{Bestimmen Sie den Wert von $m$ in der Gleichung}
    \begin{align}
    6x^3-7x^2-16x+m=0\,, \quad \overset{x=2}{\Rightarrow} \quad -12 + m \overset{!}{=} 0 \quad \Rightarrow \quad \uuline{m=12}.
    \end{align}
    \emph{wenn eine Wurzel der Gleichung den Wert 2 hat. Berechnen Sie auch die beiden anderen Wurzeln.} Wir bestimmen die restlichen Wurzeln per Polynomdivision und $pq$-Formel: 
    \begin{align}
        &\begin{array}{r@{} r@{} r@{} r@{} r}
            (6x^3 &{}-7x^2\hp{)}&{}-16x\hp{)}&{}+12) &\;:(x-2) = \uline{6x^2+5x-6} \\
          -(6x^3 &{}-12x^2&{}) \\ 
          \cmidrule{1-3}
                & 5x^2 &{} -16x &{}+12\hp{)}\\
                &-(5x^2&{} -10x&{}\hp{+12}) \\
          \cmidrule{2-4}
                & &{}-6x&{}+12\hp{)}\\
                & -&(-6x&{}+12)\\
          \cmidrule{3-4} 
                & & & 0
        \end{array}\\
        & \Rightarrow \quad 0 \overset{!}{=} x^2 +\frac{5}{6}x -1 \quad \Rightarrow x_{1/2} = -\frac{5}{12} \pm \sqrt{\frac{25}{144}+1} \quad \Rightarrow \quad \uuline{x_1 = \frac{2}{3}, x_2 = -\frac{3}{2}}.
    \end{align}
    \item \emph{Die Zahlen 2 und 3 seien Wurzeln der Gleichung}
    \begin{align}
    2x^3+mx^2-13x+n=0\,.
    \end{align}
    \emph{Bestimmen Sie die Zahlenwerte von $m$ und $n$, und geben Sie die dritte Wurzel an.}
    Zu Bestimmung der Lösung setzen wir die beiden Wurzeln in die Gleichung ein: 
    \begin{align}
        &x_1 = 2: \quad 4m + n = \hp{-}10 \quad (1)\\
        &x_2 = 3: \quad 9m + n = -15 \quad (2)\\
        & \Rightarrow \quad (1)-(2): \quad -5m = 25 \quad \Rightarrow \quad \uuline{m= -5}, \quad \uuline{n=30}.
    \end{align}
    Die dritte Wurzel ermitteln wir durch Polynomdivision mit $(x-2)(x-3) = x^2 - 5x+6$
    \begin{align}
        \begin{array}{r@{} r@{} r@{} r@{} r}
            (2x^3 &{}-5x^2\hp{)}&{}-13x\hp{)}&{}+30) &\;:(x^2-5x+6) = 2x+5 \quad \Rightarrow  \quad \uuline{x_3 = -\frac{5}{2}}\\
          -(2x^3 &{}-10x^2&{}+12x) \\ 
          \cmidrule{1-3}
                & 5x^2 &{} -25x &{}+30\hp{)}\\
                &-(5x^2&{} -25x&{}+30)\\
          \cmidrule{2-4} 
                & & & 0
        \end{array}\\
    \end{align}
\end{enumerate}

\newpage
\paragraph{Aufgabe 4: } \emph{Nullstellenraten} \hfill (Zusatzaufgabe)\\[0.2cm]
\emph{Finden Sie jeweils mindestens eine Nullstelle der folgenden Ausdrücke und spalten Sie diese als Linearfaktor $(x-x_0)$ vom Polynom ab.}
    \begin{enumerate}[label=(\alph*), labelindent=1em,labelsep=0.5cm]
        \item $x^3-5x^2+8x-4, \quad \uuline{x_0=1}$ oder $\uuline{x_0=2}$ \\
        $= (x-1)(x^2-4x+4) = (x-1)(x-2)^2$ \quad (stückweises konstruieren ohne Polynomdivision)
        \item $x^4-2x^3-13x^2+9x+9, \quad \uuline{x_0=3}$ \hfill(einzige reelle Nullstelle)\\
        $=(x+3)(x^3-5x^2+2x+3)$
        \item $x^4-3x^2+3x+2, \quad \uuline{x_0=-2}$ \hfill(einzige reelle Nullstelle) \\
        $=(x+2)(x^3-2x^2+x+1)$
        \item $x^5-x^4-3x^3+3x^2+x-1, \quad \uuline{x_0 = 1}$\\[-3mm]
        $=(x-1)(x^4-3x^2+1)$ \hfill die anderen Nullstellen sind: $x_{2/3/4/5} = \frac{\pm 1 \pm \sqrt{5}}{2}$\\
        $= (x-1)(x^2-x-1)(x^2+x-1)$
    \end{enumerate}
%
\paragraph{Aufgabe 5: } \emph{Partialbruchzerlegung} \hfill Ziel: (a) bis (c)\\[0.2cm]
Schreiben Sie, so weit möglich, als Summe von Partialbrüchen.
\begin{enumerate}[label=(\alph*)]
    \item $\frac{x-5}{x^2-2x-3} \overset{!}{=} \frac{\alpha}{x+1} + \frac{\beta}{x-3} = \frac{(x-3)\alpha + (x+1)\beta}{x^2-2x -3} = \frac{(\alpha+\beta)x +\beta - 3\alpha}{x^2-2x-3}$\\[2mm]
    Koeffizientenvergleich im Zähler: $\quad\begin{rcases}
        x^1: \quad\alpha+\beta = 1 \\ x^0: \quad3\alpha - \beta = 5
    \end{rcases}$ Addition: $\uline{\alpha = \frac{3}{2}}, \quad \uline{\beta = -\frac{1}{2}}$ \\
    $\Rightarrow \uuline{\frac{x-5}{x^2-2x-3} = \frac{3}{2(x+1)} - \frac{1}{2(x-3)}}$
    \item $\frac{x^2+1}{x^2-1} \overset{!}{=} \frac{\alpha}{x+1} + \frac{\beta}{x-1} + \gamma = \frac{(x-1)\alpha + (x+1)\beta + (x^2-1)\gamma}{x^2-1}$\\[2mm]
    Koeffizientenvergleich im Zähler: $\quad\begin{rcases}
        x^2: \quad\uline{\gamma = 1} \\ x^1: \quad\alpha + \beta = 0 \\ x^0: \quad1 = \beta - \alpha - \gamma
    \end{rcases} \quad \uline{\beta = 1}, \uline{\alpha = -1}$\\
    $\Rightarrow \uuline{\frac{x^2+1}{x^2-1} = 1 - \frac{1}{x+1}+\frac{1}{x-1}}$\\
    Alternativ kann auch eine Poylnomdivision mit Rest durchgeführt werden mit anschließender Partialbruchzerlegung des Restglieds:
    $\textstyle\frac{x^2}{x^2-1} = 1 + \frac{2}{x^2-1}.$
    \item $\frac{2x^2-3x+1}{x^3-5x^2+8x-4} \overset{!}{=} \frac{\alpha}{x-1} + \frac{\beta}{x-2} + \frac{\gamma}{(x-2)^2} = \frac{(x^2-4x+4)\alpha + (x^2-3x+2)\beta + (x-1)\gamma}{(x-1)(x-2)^2}$\\[2mm]
    Koeffizientenvergleich: $\quad\begin{rcases}
        x^2: \quad\alpha + \beta = 2 \\
        x^1: \quad-4\alpha - 3\beta + \gamma = -3 \quad (2) \\
        x^0: \quad 4\alpha +2\beta - \gamma = \hp{-}1 \quad (3)
    \end{rcases} \quad (2) + (3): \quad \beta = 2 \Rightarrow \uline{\alpha = 0}, \uline{\gamma = 3}$\\
    $\Rightarrow \uuline{\frac{2x^2 - 3x + 1}{x^3 - 5x^2 +8x-4} = \frac{2}{x-2} + \frac{3}{(x-2)^2}}$\\
    Die Faktorisierung des Nenners haben wir in Aufgabe 4a bereits gesehen. Wir müssen die Vielfachheit der Nullstelle in unserem Ansatz berücksichtigen. Durch $\alpha = 0$ sehen wir, dass die rechte Seite bei $x=1$ keine Pollstelle hat. Die Partialbruchzerlegung hat uns also den Limes $x\to 1$ der linken Seite verschafft.
    \item $\frac{2x^4-4x^3-5x^2+(\sqrt{2}-7)x+\sqrt{2}+12}{x^2-2x-3}$\\
    Da der Zähler von höherem Grade ist als das Nennerpolynom, wird zunächst eine Polynomdivision durchgeführt: 
    \begin{align}
        \begin{array}{r@{} r@{} r@{} r@{} r@{} r}
            (2x^4 &{} -4x^3 &{}-5x^2\hp{)}&{}+(\sqrt{2}-7)x&{}+\sqrt{2}+12) &\;:(x^2 - 2x-3) = \uline{2x^2+1+ \frac{(\sqrt{2}-5)x + \sqrt{2}+15}{x^2-2x-3}} \\[-3mm]
          -(-2x^4 &{}-4x^3&{}-6x^2) \\ 
          \cmidrule{1-3}
                &{} &{} x^2 &{} +(\sqrt{2}-7)x &{} +\sqrt{2} + 12\hp{)}\\
                &{} &{}-(x^2&{} -2x&{}-3) \\
          \cmidrule{3-5}
                &{} &{} &{} (\sqrt{2}-5)x &{}+\sqrt{2} + 15
        \end{array}
    \end{align}
    Die Nullstellen des Nenners lauten: $x_1 = 3, x_2 = -1$. Damit folgt 
    \begin{align}
        &\frac{(\sqrt{2}-5)x + \sqrt{2}+15}{x^2-2x-3} \overset{!}{=} \frac{\alpha}{x-3} + \frac{\beta}{x+1} = \frac{(x+1)\alpha + (x-3)\beta}{x^2-2x-3}\\
        &\text{Koeffizientenvergleich:} \quad\begin{rcases}
            x^1: \alpha + \beta = \sqrt{2} -5 \\
            x^0: \alpha - 3\beta = \sqrt{2} + 15
        \end{rcases} \quad \text{Differenz:} \quad \uline{\beta = -5}, \quad\Rightarrow \uline{\alpha = \sqrt{2}}\\
        &\Rightarrow \uuline{\frac{2x^4-4x^3-5x^2+(\sqrt{2}-7)x+\sqrt{2}+12}{x^2-2x-3} = 2x^2 + 1 + \frac{\sqrt{2}}{x-3} - \frac{5}{x+1}}
    \end{align}
\end{enumerate}
%

% \paragraph{Aufgabe 6: } \emph{Polynome in Ungleichungen}\\[0.2cm]
% Beweisen Sie die folgenden Ungleichungen.
% \begin{enumerate}[label=(\alph*)]
%     \item $(1+a+a^2)^2<3(1+a^2+a^4)\,,\qq{für} a\in\mathbb{R}\!\smallsetminus\!\{1\}$
%     \item $x^4-x^2-6x+10>0\,,\qq{für} x\in\mathbb{R}$
% \end{enumerate}

\paragraph{Aufgabe 6: } \emph{Summen}\hfill Ziel: (a) bis (b)\\[0.2cm]
Vereinfachen bzw. berechnen Sie die folgenden Summen. 

\begin{enumerate}[label=(\alph*), labelindent=1em,labelsep=0.5cm]
    \item $\underbrace{\sum_{k=1}^n \frac{1}{2k}}_{\displaystyle =\frac{1}{2}+\frac{1}{4} + \frac{1}{2n}} + \underbrace{\sum_{l=0}^{n-1} \frac{1}{2l+1}}_{\displaystyle 1+\frac{1}{3}+\frac{1}{5}+\hdots+\frac{1}{2n-1}} = 1 + \frac{1}{2} + \frac{1}{3} + \frac{1}{4} + \hdots \frac{1}{2n} = \uuline{\sum_{k=0}^{2n} \frac{1}{k}}$
    \item $\sum_{k=1}^n \frac{1}{k(k+1)}$
    \item $\sum_{k=0}^n x^k$ 
\end{enumerate}

\emph{Hinweis zu (b)}: Der Term in der Summe kann mithilfe von Partialbruchzerlegung vereinfacht und die entstehende Summe auseinandergezogen werden. 
\begin{align}
    \frac{1}{k(k+1)} = \frac{A}{k} + \frac{B}{k+1} \quad \Rightarrow \quad 1 = A(k+1) + Bk = (A+B) k + A.
\end{align}
Der Koeffizientenvergleich liefert $A=1, B=-1$ und wir können weiter vereinfachen 
\begin{align}
    \sum_{k=1}^n \frac{1}{k(k+1)} &= \sum_{k=1}^n \frac{1}{k} - \sum_{k=1}^n \frac{1}{k+1} \qq{Indexverschiebung} k+1 = l \\
    & =\sum_{k=1}^n \frac{1}{k} - \sum_{l=2}^{n+1} \frac{1}{l} = \qty(1 + \cancel{\sum_{k=2}^n}) - \qty(\cancel{\sum_{l=2}^n} + \frac{1}{n+1}) = \uuline{\frac{n}{n+1}}.
\end{align}

\emph{Hinweis zu (c)}: Hierfür kann ein expliziter Ausdruck gefunden werden, wenn man die Formel mit $(1-x)$ multipliziert und analog vorgeht wie in (b) 
\begin{align}
        (1-x) \sum_{k=0}^n x^k &= \sum_{k=0}^n x^k - x \sum_{k=0}^n x^k = \sum_{k=0}^n x^k - \sum_{k=0}^n x^{k+1} \qq{Indexverschiebung} l = k+1 \\
        &= \sum_{k=0}^n x^k - \sum_{l=1}^{n+1} x^l = 1 + \cancel{\sum_{k=1}^n x^k} - \qty(\cancel{\sum_{l=1}^n x^l} + x^{n+1}) = 1 - x^{n+1} \\
        \Rightarrow \quad \sum_{k=0}^n x^k &= \uuline{\frac{1 - x^{n+1}}{1-x}}.
\end{align}
\Titelbanner{5}{Exponentialfunktionen\\
                Logarithmen\\
                Natürliche Exponentialfunktion}

\paragraph{Aufgabe 1: } \emph{Logarithmische und Exponentialgleichungen}\\[0.2cm]
Lösen Sie jeweils für $x$ bzw. $y$. Beachten Sie mögliche Fallunterscheidungen bezüglich der auftretenden Konstanten.\\[-1.3em]
\begin{enumerate}[label=(\alph*)]
    \item $a\,2^x=\e^{bx}$
    \item $x=49^{1-\log_7(2\sqrt{x})}+5^{-\log_5(4x)}$
    \item $\e^{3ax}-2^{a+1}\e^{2ax}+4^a\e^{ax}=\log_b(1)$
    \item $\sqrt[x^2-1]{a^3}\cdot\sqrt[2x-2]{a}\cdot\sqrt[4]{a^{-1}}=1$
    \item $\log_ax+\log_ay=2\,, \quad\log_bx-\log_by=4$
    \item $3\log_{xa^2}x+\frac{1}{2}\log_\frac{x}{\sqrt{a}}x=2$
\end{enumerate}

\paragraph{Aufgabe 2: } \emph{Verdopplungszeit}\\[0.2cm]
Der Wissenszuwachs eines Physik-Studenten mit Anfangswissen $A$ sei beschrieben durch
\begin{align*}
    W(t)=A\e^{c t}\,, \hspace{1.5cm} c>0\,,
\end{align*}
sodass $W(t)$ die "`Menge"' an Wissen zur Zeit $t$ gibt.
\begin{enumerate}[label=(\alph*)]
    \item Nach welcher Zeit $\tau_2$ hat das Wissen eines Studenten auf das Doppelte zugenommen?
    \item Nach welcher Zeit $\tau_n$ hat das Wissen eines Studenten auf das $n$-Fache zugenommen?
    \item Drücken Sie $\tau_n$ als Vielfaches von $\tau_2$ aus. Welcher Zusammenhang besteht in den Fällen $n=3$ und $n=4$? Welche Aussage können Sie für $n=2^m$, $m\in\mathbb{N}$, treffen?
\end{enumerate}

\newpage
\paragraph{Aufgabe 3: } \emph{Hyperbelfunktionen}\hfill (Zusatzaufgabe)\\[0.2cm]
Es seien die Funktionen
\begin{align*}
    f(x)&:=\frac{1}{2}\qty(\e^{x}+\e^{-x})\,,\\
    g(x)&:=\frac{1}{2}\qty(\e^{x}-\e^{-x})
\end{align*}
definiert.
\begin{enumerate}[label=(\alph*)]
    \item Finden Sie jeweils einen Ausdruck für $f(2x)$ und $g(2x)$ in Abhängigkeit der Funktionen einfacher Argumente, $f(x)$ und $g(x)$.
    \item Finden Sie jeweils einen Ausdruck für $f(x+y)$ und $g(x+y)$ in Abhängigkeit von $f(x)$ und $f(y)$ sowie $g(x)$ und $g(y)$. Vergleichen Sie mit dem Ergebnis aus (a), indem Sie $x=y$ setzen.
    \item Schreiben Sie die Funktionen $f(x)$ und $g(x)$ in Reihendarstellung.
    \item Bestimmen Sie die Umkehrfunktion $g^{-1}$ von $g(x)$.
    % \item* Es sei eine dritte Funktion definiert als $h(x):=f(x)/g(x)$. Überzeugen Sie sich davon, dass diese Funktion mindestens einen \emph{Fixpunkt} besitzt, d.h. dass mindestens eine Zahl $u\in\mathbb{R}$ existiert, für die $h(u)=u$.
\end{enumerate}
% \Titelbanner{6}{Trigonometrische Funktionen\\Ebene Trigonometrie}

\paragraph{Aufgabe 1: } \emph{Additionstheoreme}\\[-1em]
\begin{enumerate}[label=(\alph*)]
\item Leiten Sie das Additionstheorem für Kosinusfunktionen aus dem für Sinusfunktionen her.
% \begin{align*}
%     \text{Es gilt:} \quad \sin(x\pm y) = \sin(x)\cos(y) \pm \cos(x) \sin(y)
% \end{align*}
\item Leiten Sie das Additionstheorem für Tangensfunktionen her,
\begin{align*}
    \tan(x\pm y)=\frac{\tan(x)\pm \tan(y)}{1\mp \tan(x)\tan(y)}\,.
\end{align*}
\item Zeigen Sie, dass für Doppelwinkelfunktionen gilt:
\begin{itemize}
    \item $\sin(2x)=2\sin(x)\cos(x)\,,$
    \item $\cos(2x)=2\cos^2\qty(x)-1\,.$
\end{itemize}
\end{enumerate}

\paragraph{Aufgabe 2: } \emph{Trigonometrische Umformungen I}\hfill Ziel: (a) bis (c)\\[0.2cm]
Zeigen Sie die Richtigkeit der folgenden Identitäten.
\begin{enumerate}[label=(\alph*)]
    \item $\dfrac{\cos(\alpha)+\sin(\alpha)}{\cos(\alpha)-\sin(\alpha)}=\tan\qty(\frac{\pi}{4}+\alpha)$
    \item $\dfrac{1+\sin (2\alpha)}{\cos (2\alpha)}=\dfrac{1+\tan(\alpha)}{1-\tan(\alpha)}=\tan\qty(\frac{\pi}{4}+\alpha)$
    \item $2\cos\qty(\frac{x+y}{2})\cos\qty(\frac{x-y}{2})=\cos(x)+\cos(y)$
    \item $\cot(\alpha)\cot(\beta)+\cot(\alpha)\cot(\gamma)+\cot(\beta)\cot(\gamma)=1 \text{ für } \alpha+\beta+\gamma=\pi$
\end{enumerate}

\paragraph{Aufgabe 3: } \emph{Trigonometrische Umformungen II}\hfill Ziel: (a) und (b)\\[0.2cm]
Formen Sie die folgenden Ausdrücke so um, dass sie sich einfach logarithmieren lassen. Das heißt, die Terme sollen möglichst in Produkte, Quotienten und Potenzen umgeformt werden.\\[-1.3em]

\begin{enumerate}[label=(\alph*)]
    \item $1+\cos(\alpha) + \cos\qty(\frac{\alpha}{2}), \qq{\emph{Hinweis:} Es ist} \cos\qty(\frac{\pi}{3})=\dfrac{1}{2}\,.$
    \item $\frac{2\sin(\beta) - \sin (2\beta)}{2\sin(\beta)+2\sin(2\beta)}$
    \item $\sin(\alpha)+\sin(\beta)+\sin(\gamma), \qq{für} \alpha+\beta+\gamma = \pi$
\end{enumerate}
%
\paragraph{Aufgabe 4: } \emph{Goniometrische Gleichungen und Gleichungssysteme}\hfill Ziel: (a) bis (c)\\[0.2cm]
Bestimmen Sie alle Lösungen der folgenden Gleichungen. Es reicht aus, die Lösungen in impliziter Form stehenzulassen!
\begin{enumerate}[label=(\alph*)]
    \item $\sin(x)+\cos(x)= 1$
    \item $\cos\qty(\frac{x+y}{2})\cos\qty(\frac{x-y}{2})=\frac{1}{2}\,, \hspace{0.5cm}\cos(x)\cos(y)=\frac{1}{4}$
    \item $\sin(3x)=\cos(2x)$
    \item $a\qty(3\cos^2(x)+\sin(x)\cos(x))-b\qty(3\sin^2(x)-\sin(x)\cos(x))=2a-b$
\end{enumerate}
\emph{Hinweis:} Versuchen Sie in (c) und (d) die Gleichung so umzuformen, dass nur noch eine Funktionsart auftritt. Dann reduziert sich das Problem zu einer Nullstellensuche eines Polynoms.

\paragraph{Aufgabe 5: } \emph{Dreiecksfläche}\\[0.2cm]
Berechnen Sie die Fläche eines Dreiecks, wenn die Seiten $a$ und $b$ sowie die Länge $w$ der Winkelhalbierenden des Winkels zwischen diesen Seiten gegeben sind.
% %
% \paragraph{Aufgabe 6: } \emph{Sehnen im Kreis}\\[0.2cm]
% Durch einen Punkt auf einem Kreis vom Radius $r$ seien zwei Sehnen der Längen $a$ und $b$ gelegt. Wenn man die Schnittpunkte der Sehnen mit der Peripherie untereinander geradlinig verbindet, erhält man ein Dreieck. Bestimmen Sie dessen Flächeninhalt $A$.

% \Titelbanner{7}{Grundlagen der Differentialrechnung\\Kurvendiskussion}

\paragraph{Aufgabe 1: } \emph{Ableitungen I}\\[0.2cm]
Berechnen Sie die Ableitungen der folgenden Funktionen.
\begin{enumerate}[label=(\alph*)]
    \item $Q(r)=\frac{3r^2}{2}\qty(\frac{1}{2}+\ln\frac{r}{r_0})$
    \item $f(x)=\cos^4(3tx)-\sin^4(3tx)$
    \item $S(\tau)=(\tau-1)\e{\tau}+\frac{\tau^2}{4}\qty(2\ln\tau-1)$
    \item $y(x)=\frac{\exp(2x)}{25}\qty[(5x-4)\sin x+(10x-3)\cos x]$
    \item $F(x)=-\frac{k}{\sqrt{(x-x_0)^2+(y-y_0)^2}}$
    \item $N(z)=\frac{2\cos\qty(\frac{z}{2})}{\sqrt{1+\cos(z)}}\qty[\ln\qty(\cos\frac{z}{4}+\sin\frac{z}{4})-\ln\qty(\cos\frac{z}{4}-\sin\frac{z}{4})]$
\end{enumerate}
%
\paragraph{Aufgabe 2: } \emph{Ableitungen II}\\[0.2cm]
Finden Sie die $n$-te Ableitung der folgenden Funktionen.
\begin{multicols}{2}
    \begin{enumerate}[label=(\alph*)]
        \item $f(x)=x^n$
        \item $f(x)=\e{kx}+\e{-kx}$
        \item $f(x)=x^{n-1}$
        \item $f(x)=a^x$ 
        \item $f(x)=\frac{1}{1-x}$
        \item[(f)*] $f(x)=\frac{x}{1-x-x^2}$ 
    \end{enumerate}
\end{multicols}
%
\paragraph{Aufgabe 3: } \emph{Kurvendiskussion I}\\[0.2cm]
Ein zweiatomiges Molekül lässt sich näherungsweise durch das sogenannte ``Morse-Potential'' beschreiben,
\begin{align*}
U(x)=D\qty(\e{-2\alpha x}-2\e{-\alpha x})\,, \hspace{1cm} D,\alpha=\operatorname{const}.
\end{align*}
\begin{itemize}
\item Bestimmen Sie Nullstellen und lokale Extrema der Funktion $U(x)$ sowie deren Verhalten für $x \to\pm\infty$.
\item Skizzieren Sie die Funktion $U(x)$ für $D=\alpha=1$ im Intervall $x\in[-1,5]$.
\end{itemize}
%
\paragraph{Aufgabe 4: } \emph{Kurvendiskussion II}\\[0.2cm]
Die Bewegung eines Teilchens mit Drehimpuls $L$ und Energie $E$ in der gekrümmten Raumzeit eines Schwarzen Loches der Masse $m$ wird beschrieben durch das Potential
\begin{align*}
U(r)=\frac{E}{2}-\frac{Em}{r}+\frac{L^2}{2r^2}-\frac{mL^2}{r^3}\,, \hspace{1cm} r>0\,.
\end{align*}
\begin{itemize}
\item Bestimmen Sie Nullstellen und lokale Extrema der Funktion $U(r)$ sowie deren Verhalten für $\to\infty$ und $\to 0$.
\item Setzen Sie $m=\frac{1}{2}$. Welche Bedingungen an $E$ und $L$ müssen erfüllt sein, damit $U(r)$ zwei, ein oder keine lokalen Extrema besitzt?
\item Skizzieren Sie die Funktion $U(r)$ für $E=1$ und $L=2$ (nicht maßstabsgerecht).
\end{itemize} 
%
\paragraph{Aufgabe 5*: } \emph{Gewöhnliche Differentialgleichungen}
\begin{enumerate}[label=(\alph*)]
\item Finden Sie eine Funktion $f(x)$, welche die folgende Gleichung erfüllt:
\begin{align*}
f''(x)=a^2f(x)+bx\,.
\end{align*}
\item Finden Sie eine Funktion $f(x)$, welche die folgende Gleichung erfüllt:
\begin{align*}
f'(x)=\qty(1+\ln(x))f(x)\,.
\end{align*}
\end{enumerate}
% \Titelbanner{8}{Die Methode der vollständigen Induktion}

\paragraph{Wiederholung des Beweisverfahrens}$~$

\begin{itemize}
    \item Zeige Richtigkeit einer Aussage für einen Startwert 
    \item Zeige: ``Aussage richtig für $k$'' impliziert ``Aussage richtig für $k+1$''
\end{itemize}

\emph{Beispiel 1:}\\
rekursive Bildungsvorschrift $a_{n+1} = a_n + 2(n+1), \quad a_1 = 1$. Wir wollen nun folgende explizite Darstellung beweisen:\footnote{Solche Ausdrücke kann Mathematica mit Hilfe des Befehls ``FindSequenceFunction$[\hdots]$'' erraten.} 
\begin{align}
    a_n = n^2 + n-1
\end{align}
\begin{enumerate}
    \item[(IA)] Induktionsanfang $n=1: a_1 = 1^2 +1 - 1 = 1 \quad\checkmark$ 
    \item[(IV)] Induktionsvoraussetzung $n=k: \quad a_k = k^2 + k -1 $
    \item[(IB)] Induktionsbehauptung $n=k+1: \quad a_{k+1}= (k+1)^2 + (k+1) - 1$
    \begin{proof}$~$\\[-1.5cm]
        \begin{align}
            \qquad a_{k+1} = a_k + 2(k+1) = k^2 +k-1 + 2(k+1) = (k+1)^2 + k.
        \end{align}
    \end{proof}
\end{enumerate}

\emph{Beispiel 2:}\\
Reihe $S_n = 1\cdot \ln(\frac{2}{1}) + 2\cdot \ln(\frac{3}{2}) + \hdots + (n-1) \cdot \ln(\frac{n}{n-1}) = \sum_{k=1}^{n-1} k \cdot \ln(\frac{k+1}{k})$
\begin{align}
        \text{Vermutung:} \quad S_n = n \cdot \ln(n) - \ln(n!)
\end{align}
\begin{enumerate}
    \item[(IA)] $n=2: S_n = 2 \cdot \ln(2) - \ln(2!) = \ln(2) \quad\checkmark$ 
    \item[(IV)] $n=k: \quad S_k = k \ln(k) - \ln(k!) $
    \item[(IB)] $n=k+1: \quad S_{k+1}= (k+1)\ln(k+1) - \ln((k+1)!)$
    \begin{proof}$~$\\[-1.5cm]
        \begin{align}
            \qquad S_{k+1} &= S_k + k \ln(\frac{k+1}{k}) \\
                           &= \cancel{k \ln(k)} - \ln(k!) + k \ln(k+1) - \cancel{k \ln(k)} \\
                           &= (k+1) \ln(k+1) \underbrace{-\ln(k+1)-\ln(k!)}_{-\ln((k+1)!)}.
        \end{align}
    \end{proof}
\end{enumerate}

\paragraph{Aufgabe 1: } \emph{Rekursive und explizite Zuordnungsvorschrift}
\begin{align}
    & a_1 = 1, a_2 = 3, a_3 = 7, a_4 = 15 \quad \Rightarrow \qq{Vermutung: } a_n=2^n-1 \hspace{0.5cm} \Rightarrow \hspace{0.5cm} a_{n+1}=2^{n+1}-1 \\
    &a_{k+1}=2a_k+1=2\qty(2^k-1)=2^{k+1}-1.
\end{align}\hfill $\Box$
%
\paragraph{Aufgabe 2: } \emph{Vollständige Induktion I} $\quad S_n = 1^2 -2^2 + 3^2 -4^2 + \hdots + (-1)^{n-1}n^2 = (-1)^{n-1}\frac{n(n+1)}{2}$
\begin{enumerate}
    \item[(IA)] $n=1: \quad S_1 = 1. \quad\checkmark$ 
    \item[(IV)] $n=k: \quad S_k=(-1)^{k-1}\frac{k(k+1)}{2}$
    \item[(IB)] $n=k+1: \quad S_{k+1}=(-1)^{k}\frac{(k+1)(k+2)}{2}$
    \begin{proof}$~$\\[-1.5cm]
        \begin{align}
            \qquad S_{k+1} &= S_k + (-1)^k (k+1)^2 =(-1)^{k-1}\frac{k(k+1)}{2}+(-1)^k(k+1)^2\\
            &=(-1)^k(k+1)\qty[\frac{(-1)^{-1}k}{2}+(k+1)]=\frac{(-1)^{k}(k+1)(k+2)}{2}.
        \end{align}
    \end{proof}
\end{enumerate}
%
\paragraph{Aufgabe 3: } \emph{Vollständige Induktion II} $\quad S_n=1^3+2^3+3^3+\dots+n^3 = \qty[\frac{n(n+1)}{2}]^2$

\begin{enumerate}
    \item[(IA)] $n=1: \quad S_1 = 1. \quad\checkmark$ 
    \item[(IV)] $n=k: \quad S_k=\qty[\frac{k(k+1)}{2}]^2$
    \item[(IB)] $n=k+1: \quad S_{k+1}=\qty[\frac{(k+1)(k+2)}{2}]^2$\\
    \begin{proof}$~$\\[-1.7cm]
        \begin{align}
            \qquad S_{k+1} &= S_k+(k+1)^3=\qty(\frac{k+1}{2})^2\qty[k^2+4(k+1)]=\qty[\frac{(k+1)(k+2)}{2}]^2.
        \end{align}
    \end{proof}
\end{enumerate}
%

\paragraph{Aufgabe 4: } \emph{Vollständige Induktion III} $\quad S_n = (n-1)^3 + n^3 + (n+1)^3, \quad S_n = 9m \qq{mit} m \in \mathbb{N}$

\begin{enumerate}
    \item[(IA)] $n=1: \quad S_1 = 0 + 1^3 + 2^3 = 9. \quad\checkmark$ 
    \item[(IV)] $n=k: \quad S_k = (k-1)^3 + k^3 + (k+1)^3, \quad S_k = 9m \qq{mit} m \in \mathbb{N}$
    \item[(IB)] $n=k+1: \quad S_{k+1} = k^3 + (k+1)^3 + (k+2)^3, \quad S_k = 9m' \qq{mit} m' \in \mathbb{N}$\\
    \begin{proof}$~$\\[-1.4cm]
        \begin{align}
            \qquad S_{k+1}&=k^3+(k+1)^3+(k+2)^3=\underbrace{(k-1)^3+k^3+(k+1)^3}_{S_n}+(k+2)^3-(k-1)^3\\
            &=9m+(k+2)^3-(k-1)^3=9m+9(k^2+k+1)=9m' \text{ mit } m'=m+k^2+k+1\in\mathbb{N}
        \end{align}
    \end{proof}
\end{enumerate}
%
\paragraph{Aufgabe 5: } \emph{Die Suche nach der richtigen Summenformel}\\
\begin{align}
\text{Vermutung: }\hspace{0.2cm} S_n=\frac{(1-x)(2-x)\dots(n-x)}{n!} \hspace{0.5cm}
\end{align}

\begin{enumerate}
    \setlength{\mathindent}{0cm}
    \item[(IA)] $n=1: \quad S_1 = 1-x. \quad\checkmark$ 
    \item[(IV)] $n=k: \quad S_k = \frac{(1-x)(2-x)\dots(k-x)}{k!}$
    \item[(IB)] $n=k+1: \quad S_{k+1} = \frac{(1-x)(2-x)\dots((k-x)(k+1-x)}{(k+1)!}$\\
    \begin{proof}
        \begin{align}
            S_{k+1}&=S_k+(-1)^{k+1}\frac{x(x-1)(x-2)\dots(x-k+1)(x-k)}{(k+1)!}\\
            &=\frac{(1-x)(2-x)\dots(k-x)}{k!}+(-1)^{k+1}\frac{\textcolor{red!70!black}{x(x-1)(x-2)\dots(x-k+1)(x-k)}}{(k+1)!}\\
            &=\frac{1}{(k+1)!}\qty[(1-x)(2-x)\dots(k-x)(k+1)+(-1)^{k+1}\textcolor{red!70!black}{(-1)^k x(1-x)(2-x)\dots(k-x-1)(k-x)}]\\
            &=\frac{(1-x)(2-x)\dots(k-x)(k-x+1)}{(k+1)!}\underbrace{\qty[\frac{k+1}{k-x+1}+(-1)^{2k+1}\frac{x}{k-x+1}]}_1\\
            &=\frac{(1-x)(2-x)\dots(k-x)(k-x+1)}{(k+1)!}
        \end{align}
    \end{proof}
\end{enumerate}

\paragraph{Aufgabe 6: } \emph{Fibonacci-Zahlen} $\quad a_n=\frac{1}{\sqrt{5}}\qty(x_+^n-x_-^n) \qq{mit}x_\pm = \frac{1\pm \sqrt{5}}{2}$\hfill (Zusatzaufgabe)

\begin{enumerate}
    \setlength{\mathindent}{0cm}
    \item[(IA)] $n=0: \quad a_0 = 0, a_1 = 1. \quad\checkmark$ 
    \item[(IV)] $n=k: \quad a_k = \frac{1}{\sqrt{5}}\qty(x_+^k-x_-^k)$
    \item[(IB)] $n=k+1: \quad a_{k+1} = \frac{1}{\sqrt{5}}\qty(x_+^{k+1}-x_-^{k+1})$
    \begin{proof}
        \begin{align}
            a_{n+1}&=a_n+a_{n-1}=\frac{1}{\sqrt{5}}\qty(x_+^n-x_-^n+x_+^{n-1}-x_-^{n-1})\\
            &=\frac{1}{\sqrt{5}}\qty[\qty(\frac{1+\sqrt{5}}{2})^n+\qty(\frac{1+\sqrt{5}}{2})^{n-1}-\qty(\frac{1-\sqrt{5}}{2})^n-\qty(\frac{1-\sqrt{5}}{2})^{n-1}]\\
            &=\frac{1}{\sqrt{5}}\Bigg[\underbrace{\qty(\frac{1+\sqrt{5}}{2}+1)}_{\qty(\frac{1+\sqrt{5}}{2})^2}\qty(\frac{1+\sqrt{5}}{2})^{n-1}-\underbrace{\qty(\frac{1-\sqrt{5}}{2}+1)}_{\qty(\frac{1-\sqrt{5}}{2})^2}\qty(\frac{1-\sqrt{5}}{2})^{n-1}\Bigg]
            =\frac{1}{\sqrt{5}}\qty(x_+^{n+1}-x_-^{n+1})
        \end{align}
    \end{proof}
\end{enumerate}
% \Titelbanner{9}{Arithmetische \& geometrische Reihen\\
                Der binomische Satz}

\paragraph{Aufgabe 1: } \emph{Arithmetische Reihe}\\[0.2cm]
Eine Spirale bestehe aus zwei Scharen konzentrischer Halbkreise um die Punkte $A$ und $B$. Es sei $r$ der Radius des innersten Halbkreises und die Strecke $\overline{AB}=e$.
\begin{enumerate}[label=(\alph*)]\setlength{\itemsep}{-0.5ex}
\item Wie lang ist der $n$-te Halbbogen?
\item Wie lang ist der Gesamtbogen der Spirale bis dahin?
\end{enumerate}
\begin{figure}[htp]
    \centering
    \begin{tikzpicture}[scale=0.5]
        \draw[thick] (-6,0) -- (6,0);
        \draw[thick] (2.3, 0) arc(0:-180:2.3); 
        \draw[thick] (4.3, 0) arc(0:-180:4.3);
        \draw (0,-.1)node[below]{$B$} -- (0,.1); 
        \draw[{latex}-{latex}] (0,-1) --node[below]{$e$} (1.3,-1);
        \begin{scope}[shift={(1,0)}]
            \node (A) at (0.65,0.3){$r$};
            \draw (0,-.1)node[below]{$A$} -- (0,.1);
            \draw[thick] (1.3, 0) arc(0:180:1.3); 
            \draw[thick] (3.3, 0) arc(0:180:3.3); 
        \end{scope}
    \end{tikzpicture}
    \vspace{-0.5cm}
\end{figure}
%
\paragraph{Aufgabe 2: } \emph{Geometrische Folge}\\[0.2cm]
Beim Durchgang durch eine Glasplatte verliert ein Lichtstrahl $\textstyle\frac{1}{12}$ seiner Lichtstärke. Wie viele Platten muss er durchdringen, wenn er nur noch die Hälfte der ursprünglichen Lichtstärke besitzen soll?\\[0.2cm]
\emph{Hinweis:} Es genügt die Angabe des Ergebnisses in impliziter Form.
%
\paragraph{Aufgabe 3: } \emph{Geometrische Reihe}\\[0.2cm]
Es sei eine Anordnung aus zwei Glasplatten gegeben, in die ein Laserstrahl der Intensität $\mathcal{I}_\text{in}$ eingeschossen werde. Jede Platte besitze einen Transmissionskoeffizienten $T$ ($0<T<1$), d.h. dass an jeder Platte von einem Lichtstrahl der Intensität $\mathcal{I}$ ein Anteil $T \mathcal{I}$ transmittiert und ein Anteil $(1-T)\mathcal{I}$ reflektiert wird.
Bestimmen Sie die Intensität $\mathcal{I}_\text{out}$, mit der das Licht auf der anderen Seite der Anordnung austritt.\\
\begin{center}
    \vspace{-1cm}
\begin{tikzpicture}[scale=0.8]
    \node [] at(0.2,0){$\mathcal{I}_\text{out}$};
    \draw [<-,line width=0.8pt] (0.7,0)--(2.7,0);
    \draw [line width=1pt] (3.0,-2)--(3.0,2);
    \draw [line width=1pt] (6.0,-2)--(6.0,2);
    \draw [<-,line width=0.8pt] (3.3,0.15)--(5.7,0.15);
    \draw [->,line width=0.8pt] (3.3,-0.15)--(5.7,-0.15);
    \draw [<-,line width=0.8pt] (6.3,0)--(8.3,0);
    \node [] at(8.8,0){$\mathcal{I}_\text{in}$};
    \node [below] at(3,-2){$T$};
    \node [below] at(6,-2){$T$};
\end{tikzpicture}
\vspace{-2cm}
\end{center}
%
\newpage
\paragraph{Aufgabe 4: } \emph{Der binomische Satz}
\begin{enumerate}[label=(\alph*)]
\item Berechnen Sie $(1+a)^6+(1-a)^6$.
\item Schreiben Sie die allgemeine Binomialformel für $(1+x)^n$ und $(1-x)^n$ auf und setzen Sie anschließend $x=1$. In letzterem Falle sei $n\ge 1$.
\item Es seien die Funktionen $f(x)=\left(\e{x}+\e{-x}\right)/2$ und $g(x)=\left(\e{x}-\e{-x}\right)/2$ definiert. Finden Sie einen Ausdruck, der $f(nx)+g(nx)$, mit $n\in\mathbb{N}$, auf (Potenzen von) $f(x)$ und $g(x)$ zurückführt.
\end{enumerate}
%
\paragraph{Aufgabe 5: } \emph{Erzeugende Funktion}\\[0.2cm]
Betrachten Sie die Rekursionsgleichung $a_{n+1}=2a_n+1$ mit Anfangswert $a_0=0$ aus Aufgabe 1, Thema 8. Es sei eine (unbekannte) Funktion definiert als
\begin{align*}
f(x)=\sum\limits_{n=0}^\infty a_n x^n\,.
\end{align*}
\begin{enumerate}[label=(\alph*)]
\item Multiplizieren Sie beide Seiten der Rekursionsgleichung mit $x^n$, summieren Sie über $n$ (von Null bis Unendlich) ab und versuchen Sie alle auftretenden Terme durch $f(x)$ auszudrücken. Lösen Sie für $f(x)$.
\item Nutzen Sie Ihr Wissen über geometrische Reihen, um $f(x)$ wieder in Reihendarstellung zu überführen und lesen Sie die Koeffizienten vor $x^n$ ab.
\end{enumerate}
%
\paragraph{Aufgabe 6*: } \emph{Zustandssumme}\\[0.2cm]
Vereinfachen Sie den Ausdruck
\begin{align*}
\Omega_N=\sum\limits_{n_1=0}^N\sum\limits_{n_2=0}^N\sum\limits_{n_3=0}^N\dots\sum\limits_{n_k=0}^N f(n_1)f(n_2)f(n_3)\dots f(n_k)\,, \qq{wobei}f(n)=\e{-n}.
\end{align*}
Bestimmen Sie anschließend den Grenzwert $\Omega=\lim\limits_{N\rightarrow\infty}(\Omega_N)$.
%
\paragraph{Aufgabe 7*: } \emph{Fibonacci-Zahlen II}\\[0.2cm]
Wenden Sie das Verfahren aus Aufgabe 5 auf die Folge $a_{n+1}=a_n+a_{n-1}$ der Fibonacci-Zahlen mit Startwerten $a_0=0$, $a_1=1$ an. Beachten Sie, dass hier $n\ge 1$ gelten muss.
%
% \Titelbanner{10}{Rechnen mit Vektoren und Matrizen}

\paragraph{Aufgabe 1: } \emph{Linear (un)abhängige Vektoren}
\begin{enumerate}[label=(\alph*)]
\item Sind die folgenden Vektoren linear unabhängig?
\begin{align*}
\mqty( 3\\1\\-5 ),\quad \mqty( 1\\2\\0 ),\quad \mqty(-2\\0\\1)
\end{align*}
\item Für welche Werte von $t$ befinden sich die folgenden Vektoren in einer Ebene?
\begin{align*}
\mqty( 1\\4\\0 ),\quad \mqty( 3\\t\\1 ),\quad \mqty( 4\\4\\2 )
\end{align*}
\end{enumerate}
%
\paragraph{Aufgabe 2: } \emph{Skalarprodukt und Vektorprodukt}
\begin{enumerate}[label=(\alph*)]
\item Berechnen Sie die folgenden Skalarprodukte.
\begin{align*}
\mqty(1\\5\\7)\cdot\mqty(-2\\8\\3), &&
\mqty(1\\0\\ \sqrt{2})\cdot\mqty(0\\-1\\ \sqrt{2}), &&
\mqty(a\\3\\-1)\cdot\mqty(5\\2\\5a+6)
\end{align*}
\item Berechnen Sie die folgenden Vektorprodukte.
\begin{align*}
\mqty(1\\0\\0 )\times\mqty(0\\1\\0),&&
\mqty(\cos(\phi)\\ \sin(\phi) \\0)\times\mqty(-\sin(\phi)\\ \cos(\phi)\\0),&&
\mqty(1\\3\\-2)\times\mqty(1\\1\\5)
\end{align*}
\item Nutzen Sie das Skalarprodukt, um den Kosinussatz im allgemeinen Dreieck herzuleiten.
\item Nutzen Sie das Vektorprodukt, um den Sinussatz im allgemeinen Dreieck herzuleiten.
\end{enumerate}
%
\paragraph{Aufgabe 3: } \emph{Matrizen}
\begin{enumerate}[label=(\alph*)]
\item Bilden Sie -- falls möglich -- die Inversen der folgenden Matrizen mit Hilfe des Gauß-Jordan-Algorithmus.
\begin{align*}
 A =\mqty(1 & 2\\ 3 & 4),\quad B=\mqty(0 & 2\\ 3 & 0),\quad  C=\mqty(-2 & 4\\ 3 & -6)
\end{align*}
\end{enumerate}
\newpage
\begin{enumerate}[label=(\alph*),resume]
\item Matrixmultiplikation ist im Allgemeinen nicht kommutativ. Berechnen Sie $ A \cdot B- B\cdot A$.
\item Welche Bedingungen müssen die Einträge $m_i$ einer Matrix
\begin{align*}
M = \mqty(m_1 & m_2\\ m_3 & m_4)
\end{align*}
erfüllen, damit $| M\cdot \vec{v}|=|\vec{v}|$ für einen beliebigen zweikomponentigen Vektor $\vec{v}$ gilt? 
\end{enumerate}
%
%\paragraph{Aufgabe 4: } \emph{Drehungen}\\[0.2cm]
%Gegeben seien eine Matrix $\hat{R}(\phi)$ und ein Vektor $\vec{\alpha}$ als
%\begin{align*}
%\hat{R}=\begin{pmatrix*}[c]\cos(\phi)	& -\sin(\phi)	& 0\\ \sin(\phi) & \cos(\phi) & 0\\ 0 & 0 & 1\end{pmatrix*}, \quad \vec{\alpha}=\mqty(x\\y\\z).
%\end{align*}
%Berechnen Sie $\hat{R}(\phi)\cdot\vec{\alpha}$ und $\det\hat{R}(\phi)$. Welche geometrische Bedeutung hat der Fall $\phi=\frac{\pi}{2}$, $x=1$, $y=z=0$?\\[1cm]
%
\paragraph{Aufgabe 4: } \emph{Laplace'scher Entwicklungssatz}\\[0.2cm]
Berechnen Sie die Determinante der folgenden Matrix. Für welche Werte von $a$ ist die Matrix nicht invertierbar?
\begin{align*}
A=\mqty(a &3 &1 &a\\ 0 &-2 &0 &1\\ 3 &2 &-4 &-1\\ -a& 0 &-2 &0)
\end{align*}
%
\paragraph{Aufgabe 5: } \emph{Vektorprodukt}\\[0.2cm]
Das Vektorprodukt $\vec{\omega}\times \vec{u}$ zweier Vektoren
$\quad \vec{\omega}=\mqty(\omega_1\\\omega_2\\\omega_3),\quad\vec{u}=\mqty(u_1\\u_2\\u_3)$

lässt sich auch als Matrixmultiplikation $\Omega \cdot \vec{u}$ mit einer ($3\times 3$)-Matrix $\Omega$ mit den Einträgen $\pm \omega_i$, $i\in\{1,2,3\}$ schreiben. Wie sieht  $\Omega$ aus? Welche Eigenschaften hat es (Determinante, Spur, Symmetrie)?
%
%
\paragraph{Aufgabe 6*: } \emph{Höhere Dimensionen}\\[0.2cm]
Im dreidimensionalen euklidischen Raum $\mathbb{R}^3$ kann jede beliebige Drehung in Drehungen um die $x$-, die $y$- und die $z$-Achse zerlegt werden, d.h. man benötigt genau drei linear unabhängige Drehmatrizen für eine solche Zerlegung. Wie viele linear unabhängige Drehmatrizen benötigt man im $d$-dimensionalen euklidischen Raum $\mathbb{R}^d$\,?
%
%
\paragraph{Aufgabe 7*: } \emph{Diskreter Laplace-Operator}\\[0.2cm]
Gegeben sei die ($n\times n$)-Matrix $\Delta_n$, welche die Einträge -2 auf der Haupt- und 1 auf den beiden Nebendiagonalen hat,
\begin{align*}
\Delta_n=\mqty(
-2 & 1 & 0 & 0 & 0 &\dots & 0\\
1 & -2 & 1 & 0& 0 &\dots & 0\\
0 & 1 & -2 & 1 & 0 &\dots & 0\\
\vdots &\vdots &\vdots &\vdots &\vdots & &\vdots\\
0 & 0 & 0 & 0 & 0 &\dots & -2
)\,.
\end{align*}
Bestimmen Sie die Determinante von $\Delta_n$\,, indem Sie wie folgt vorgehen:
\begin{enumerate}
\item Finden Sie eine Rekursionsrelation, die $\det(\Delta_n)$ auf $\det(\Delta_{n-1})$ und $\det(\Delta_{n-2})$ zurückführt.
\item Raten Sie einen allgemeinen Ausdruck für $\det(\Delta_{n})$ und beweisen Sie ihn mittels vollständiger Induktion. Alternativ kann der Versuch unternommen werden, das Verfahren aus Thema 9, Aufgabe 5 anzuwenden.
\end{enumerate}

\end{document}